\section{Cumulative gravitational endpoint resource}
\label{sec:resource}

Let a compact elastic body have equilibrium density $\rho(\bm x)$ and center of mass at the origin. Expand a small displacement field as
\begin{equation}
\bm u(\bm x,t)=\sum_n \xi_n(t)\bm w_n(\bm x),
\end{equation}
with mass inner product
\begin{equation}
\langle u,v\rangle_\rho=\int_V\rho\,\bm u\cdot\bm v\,\dd^3x,
\qquad
\langle w_n,w_m\rangle_\rho=\mu_n\delta_{nm}.
\end{equation}
The long-wavelength tidal coupling and elastic-mode treatment are standard gravitational-antenna theory \cite{Hirakawa1976,Lobo1995}.

To first order in displacement, define the STF modal quadrupole
\begin{equation}
q_{n,ij}
=\int_V\rho\left(
 w_{n,i}x_j+w_{n,j}x_i
-\frac23\delta_{ij}\bm w_n\cdot\bm x
\right)\dd^3x .
\label{eq:qn}
\end{equation}
For a harmonic mode of physical frequency $\omega_n$, the standard quadrupole power formula gives the gravitational energy-decay rate
\begin{equation}
\boxed{
\kappa_{g,n}
=\frac{G\omega_n^4}{5c^5}
\frac{q_n:q_n}{\mu_n} .
}
\label{eq:kappan}
\end{equation}

The cumulative resource follows from one elementary completeness bound. Introduce the vector influence fields
\begin{equation}
(g^{ij})_k
=\delta_{ik}x_j+\delta_{jk}x_i
-\frac23\delta_{ij}x_k,
\end{equation}
so that $q_{n,ij}=\langle w_n,g^{ij}\rangle_\rho$. Direct contraction gives
\begin{equation}
\boxed{
\sum_{i,j,k}|(g^{ij})_k|^2=\frac{20}{3}r^2 .
}
\label{eq:20over3}
\end{equation}
For the normalized modes $\phi_n=w_n/\sqrt{\mu_n}$, Bessel's inequality applied componentwise yields
\begin{align}
\sum_n\frac{q_n:q_n}{\mu_n}
&\le
\int_V\rho\sum_{i,j,k}|(g^{ij})_k|^2\dd^3x \\
&=\frac{20}{3}I_2 .
\end{align}
Therefore
\begin{equation}
\boxed{
\sum_n\frac{q_n:q_n}{\mu_n}\le\frac{20}{3}I_2 .
}
\label{eq:qsum}
\end{equation}

If the retained modal sector satisfies $\omega_n\le\Omega$, Eqs.~\eqref{eq:kappan} and \eqref{eq:qsum} imply
\begin{equation}
\boxed{
\sum_n\kappa_{g,n}
\le
\frac{4G}{3c^5}I_2\Omega^4 .
}
\label{eq:kappasum}
\end{equation}
In an energy-normalized Markov modal representation this is
\begin{equation}
\Tr(K_g^\dagger K_g)=\sum_n\kappa_{g,n}.
\end{equation}
The trace is invariant under unitary internal modal mixing. For a countable retained modal sector, Eq.~\eqref{eq:kappasum} also proves that $K_g$ is Hilbert--Schmidt.

The theorem uses the retained carrier-scale sector of Eq.~\eqref{eq:retainedsector}, for which $\Omega=\omega_0[1+O(B/\omega_0)]$. Therefore
\begin{equation}
\Tr(K_g^\dagger K_g)
\lesssim
\frac{4G}{3c^5}I_2\omega_0^4 .
\label{eq:narrowresource}
\end{equation}
This replacement is not a statement about arbitrary endpoint modes at all physical frequencies. In particular, modes with $\omega_n\gg\omega_0$ can have off-resonant tails inside the measured envelope band while carrying gravitational rates proportional to their own $\omega_n^4$. Such sectors are outside the simple carrier-scale closure unless their contribution is separately bounded.

Neither the existence of an integrated gravitational response nor the general modal/sum-rule philosophy is new. Paik and Wagoner calculated resonant-mass absorption cross sections \cite{PaikWagoner1976}; Aguiar's review explicitly discusses the integrated cross section and notes that adding a resonant transducer produces two normal modes while retaining a finite mass/material scaling \cite{Aguiar2011}; and gravitational material-response dispersion and sum rules have been derived microscopically \cite{Srivastava2003}. Accordingly, Eqs.~\eqref{eq:20over3}--\eqref{eq:kappasum} are used here as proof lemmas, not as standalone novelty claims.
